%% 自動生成（gen.py）
\bktocchap{8-0}{空間のどこにも矢印がある}{天気図の風も、磁石のまわりも矢印でいっぱい}{6}{3}{3}{8}{J}
\bktocpart{第1部}{たくさんの変数の微分}{28}
\bktocchap{8-1}{地図と等高線：2つ以上の変数の関数}{場所を決めると、高さが1つ決まる}{1}{4}{4}{30}{C}
\bktocchap{8-2}{偏微分：ほかを止めて1つだけ動かす}{東へ進むときの坂と、北へ進むときの坂}{1}{4}{4}{64}{C}
\bktocchap{8-3}{全微分：少し動いたときの変化}{ななめに少し進むと、高さはどれだけ変わる？}{1}{5}{5}{100}{C}
\bktocchap{8-4}{いちばん急な向き：勾配 grad}{ボールが転がる向きは、等高線と直角}{1}{5}{5}{138}{C}
\bktocpart{第2部}{ベクトル場と積分}{174}
\bktocchap{8-5}{ベクトル場}{空間の点ごとに、矢印を1本ずつ}{6}{4}{4}{176}{C}
\bktocchap{8-6}{面や立体でたし合わせる：重積分}{厚さがちがう板の重さを求める}{2}{5}{5}{202}{C}
\bktocchap{8-7}{道にそってたす：線積分}{遠回りすると、仕事は変わる？}{2}{5}{5}{252}{C}
\bktocchap{8-8}{面を通りぬける量：面積分と流束}{あみの向きで、通りぬける水の量が変わる}{2}{5}{5}{292}{C}
\bktocpart{第3部}{発散と回転}{344}
\bktocchap{8-9}{わき出しと吸いこみ：発散 div}{小さな箱から、どれだけ出ていく？}{6}{5}{5}{346}{C}
\bktocchap{8-10}{うず：回転 rot}{小さな風車を置くと、回るかな？}{6}{5}{5}{392}{C}
\bktocchap{8-11}{中と外をつなぐ：ガウスの定理とストークスの定理}{中のわき出しを全部たすと、外へ出る量になる}{2}{5}{5}{436}{C}
\bktocchap{8-12}{ならす働き：ラプラシアン}{まわりの平均との差が、変化を決める}{1}{5}{5}{488}{C}
\bktocpart{第4部}{偏微分方程式}{526}
\bktocchap{8-13}{熱と波の伝わり方：偏微分方程式}{金属の棒に熱が広がるようすを、表計算で再現する}{1}{5}{5}{528}{C}
\bktocchap{8-14}{div と rot で書かれた式を読もう}{電気と磁気の4つの式を、「わき出し」と「うず」で読む}{6}{5}{5}{568}{R}
\par\vspace{2mm}
\bktocbacklist{\bkbi{こたえ}{586}\bkbi{\R{記号}{き|ごう}カード}{602}\bkbi{よく\R{使}{つか}う\R{定数}{てい|すう}・\R{単位}{たん|い}・\R{公式}{こう|しき}}{608}\bkbi{できるようになったことチェック\R{表}{ひょう}}{624}\bkbi{つぎの\R{巻}{かん}へ}{627}\bkbi{さくいん}{629}}
